\chapter{补充结果与正确性检查}
\label{app:supplementary}
本附录记录支持第3--8章的参数覆盖、逐阶段验证及次要结果。
各测量家族保留各自版本与计时边界，实验协议见附录\ref{app:reproducibility}。

\section{参数覆盖}
\label{sec:b-grids}
E1比较有1,590个有效实测配置，覆盖见\cref{tab:b-e1-counts}。
下列集合表示观测值，不假定其笛卡尔积全部测量过。

\par\medskip\noindent\begin{minipage}{\linewidth}
\centering\captionsetup{hypcap=false}
\captionof{table}[E1 精确配置数量]{包含覆盖扩展的E1精确配置数。计数为参数点，不是独立索引构建。}
\label{tab:b-e1-counts}
\footnotesize\setlength{\tabcolsep}{4pt}
\begin{tabular}{@{}lrrrrr@{}}
\toprule
方法 & Vogue & arXiv & OpenAI-1536 & OpenAI-3072 & 总计 \\
\midrule
JHQ-GPU v69 & 27 & 29 & 37 & 39 & 132 \\
CAGRA fp32 & 64 & 64 & 64 & 64 & 256 \\
CAGRA int8 & 64 & 64 & 64 & 64 & 256 \\
IVF-RaBitQ & 132 & 132 & 264 & 198 & 726 \\
IVF-PQ & 54 & 54 & 54 & 54 & 216 \\
GPU brute force & 1 & 1 & 1 & 1 & 4 \\
\midrule
总计 & 342 & 344 & 484 & 420 & 1590 \\
\bottomrule
\end{tabular}

\end{minipage}\par\medskip

\paragraph{JHQ-GPU。}
E1均用$B=B_r=8$、$M=d/8$、$k=10$和batch 1,000；Vogue为4,096个IVF列表，
其他三个为8,192。Vogue观测探针16、32、64、96、128、192、256、384、512，
$\alpha\in\{16,32,64,128\}$，共27个实际组合。
arXiv再包含探针1,024，$\alpha\in\{25,50,100,200\}$，共29个组合。
两个OpenAI均用探针32、64、128、256、512、1,024；倍率集合分别为
$\{4,8,16,32,64,128,256,512\}$和$\{2,4,8,16,32,64,128,256\}$，共37／39个组合。
所报调度路径在Vogue用四组，其余用单block。
历史v69的OpenAI旋转／质心GEMM为TF32；v72 FP32敏感性另成家族
（\cref{sec:7-precision}）。

\paragraph{CAGRA fp32与int8。}
每种表示每数据集64点：图度数32／64、中间度数128、\texttt{ivf\_pq}构建，
\texttt{itopk\_size}为32、64、96、128、192、256、384、512，搜索宽度1、2、4、8。
int8数据缩放记录quantile 0.99，主机搜索边界包含查询量化。

\paragraph{IVF-PQ。}
八位子码在768维使用192、384、768字节编码，OpenAI使用384、768、1,536字节。
每种长度测试fp16／fp32 LUT及探针16、32、64、96、128、192、256、384、512，
每数据集54点。记录的精排比例为一，列表数与JHQ相同；编码长度和显存预算不同。

\paragraph{IVF-RaBitQ与暴力搜索。}
RaBitQ包含每维4、8、9位、探针16、32、64、96、128、192、256、384、512、768、1,024。
768维与OpenAI-1536含LUT16、LUT32、QUANT4、QUANT8，OpenAI-3072含
LUT32、QUANT4、QUANT8；另用4,096列表OpenAI索引扩展默认8,192列表覆盖。
精确组合见参数记录。每数据集还含一个FP32平方L2暴力搜索点。
缺少某种模式或更高recall点不等于实测失败。

\paragraph{辅助网格。}
LUT实验固定$\alpha=100$，六数据集探针8、32、128、512，
$\glut\in\{1,2,4,8\}$，共24单元、96计时臂。
预算实验在相同六数据集用探针128／512、$A=(200,100,64,32,16,8,4,2)$，
每个$S\in\{32,64,128\}$抽取64次、$\epsilon=1$；FULL另用全部500校准查询。
留出为500条，Stella为499条。直接复用为四单元、四策略、五轮，每轮500个重复留出批次。
边界见\cref{sec:a-budget-contract,sec:a-supplement}。

\section{正确性与数值证据}
\label{sec:b-correctness}
\Cref{tab:b-checks}区分数值、算法与内存证据。
第五章排序论证针对满足所列契约的计算分数；原始向量recall另行评测。

\par\medskip\noindent\begin{minipage}{\linewidth}
\centering\captionsetup{hypcap=false}
\captionof{table}[校验覆盖范围]{验证覆盖及实际核对对象。}
\label{tab:b-checks}
\footnotesize\setlength{\tabcolsep}{4pt}
\begin{tabularx}{\linewidth}{@{}>{\raggedright\arraybackslash}p{.22\linewidth}>{\raggedright\arraybackslash}X@{}}
\toprule
边界 & 已有证据与范围 \\
\midrule
表示 & 四库全行offset／ID审计，原始输入采样FP64重编码、存储矩阵几何，已识别TF32偏差 \\
Cartesian LUT & 实数因子化；96 个计时臂；T4 覆盖四种表形式的逐候选 FP32 误差与排名诊断 \\
全局选择 & T4：270 配置、69,120 查询实例；55,296 个分组实例与严格 CPU 完整键排序的顺序一致 \\
共享界 & 48组配对；提取标量64位PTX及toolkit SASS支持完整key；未捕获驱动实际SASS \\
调度 & 250个变体／配置及234,320次查询比较，区分ID集合与并列顺序 \\
探针选择 & 1,792合成用例零失败；真实查询ID／offset检查，针对已计算路由分数 \\
预算搜索 & 2,304采样曲线单调、12个FULL枚举一致；18,432网格计数可由全池子集和复现 \\
\bottomrule
\end{tabularx}
\end{minipage}\par\medskip

较小表示审计库在按ID关联的一遍／两遍构建中主码和残差字节相同，correction差至多
$8.94\times10^{-8}$，列表分配有1、1、40行变化；Stella验证自然两遍路径
（\cref{sec:4-representation-audit}）。T4 已提供计算分数选择的严格 CPU 参考
（\cref{sec:a-t4-t5}）；同模型 Vogue 跨进程受控诊断仍未执行：归档进程即使
probe／primary dump 一致，最终 recall 仍有 0.9642／0.9644。
进程内选择检查并未定位这项独立异常。

\section{次要结果与证据状态}
\label{sec:b-secondary}
\paragraph{不利运行区间。}
调度器最大记录regret为batch 256的$1.217\times$；Vogue E1中调度／单block吞吐
为0.895--1.175。使0.97 recall领先的偏斜规则也会在低探针单元增加开销
（\cref{sec:7-6}），这些不利结果与内部调度收益一起保留。T5 跨系统比较中，
recall 0.97、batch 8 时四个数据集均由 CAGRA fp32 领先；batch 32 的边界依赖
recall 估计（\cref{sec:7-smallbatch}）。

\paragraph{表示替代方案。}
Mailman式GPU评分器为独立kernel测量，端到端替换只是乐观模型。
面向排序的主研究为FP64离线回放；另一个median编码GPU后续实验在记录目标上，
arXiv的batch 8为$1.402\times$、batch 1,000为$1.208\times$，
Vogue未达严格目标，OpenAI-3072接近平手。阈值、级值和残差训练同时改变，
不能只归因于阈值（\cref{sec:8-2}）。

\paragraph{校准与精度。}
有限池预测在查询计数稳定的条件下解释抽样风险，GT损失映射使用留出固定臂质量。
直接复用测量四个重复池的累计成本；精度干预在同一TF32训练缓存上改变add／查询GEMM。
完整FP32重训、未来查询漂移、跨设备排名及在线更新不属于这些实验。

\section{构建容量与精度锚点}
\label{sec:b-full-review}

\Cref{tab:b-scratch}区分输入暂存与独立分配的乘积数组。源码默认三个缓冲区、
配置 add 行数 65,536、每个输入／旋转浮点数组 128\,MiB。
这些是分配容量，不是当前子批大小或进程整体峰值。

\par\medskip\noindent\begin{minipage}{\linewidth}
\centering\captionsetup{hypcap=false}
\captionof{table}[独立构建缓冲容量]{默认构建缓冲容量。
末列为 $4n_b n_{\rm list}a_{\rm assign}/2^{30}$；用 $a_{\rm in}$ 代入会得到错误的一般模型。}
\label{tab:b-scratch}
\footnotesize\begin{tabular}{@{}rrrrr@{}}
\toprule
$d$ & $n_{\rm list}$ & $a_{\rm in}$ & $a_{\rm assign}$ & 分配乘积 GiB \\
\midrule
768 & 4,096 & 43,690 & 8,192 & 0.375 \\
1,024 & 32,768 & 32,768 & 32,768 & 12.000 \\
1,536 & 8,192 & 21,845 & 32,768 & 3.000 \\
3,072 & 8,192 & 10,922 & 32,768 & 3.000 \\
\bottomrule
\end{tabular}

\end{minipage}\par\medskip

\subsection{同期精度锚点}
\label{sec:b-precision-anchors}

T1 晚于 D，在同一会话重测 recall 0.97 附近的 JHQ 和竞争基线，启用 E1 工作区池。
每个 JHQ 配置在基线前后各重复一次，每次五轮计时；每配置 QPS 对两次轮中位数
取几何平均。每个 OpenAI 数据集有十二个配对配置，FP32/TF32 几何平均为
0.956／0.964。两种变体复用 TF32 训练缓存，改变 add／查询算术。

\par\medskip\noindent\begin{minipage}{\linewidth}
\centering\captionsetup{hypcap=false}
\captionof{table}[同会话精度锚点]{T1，recall 0.97、batch 1,000。
基线列依次为插值／严格实测阈值，FP32/TF32 列为插值。
这里全部方法均为同期测量，与 D 的跨会话基线比值分开。}
\label{tab:b-precision-anchors}
\footnotesize\setlength{\tabcolsep}{3pt}
\begin{tabular}{@{}lrrr@{}}
\toprule
数据集 & 0.97 处 FP32/TF32 & FP32 / RaBitQ & FP32 / CAGRA \\
\midrule
OpenAI-1536 & 0.980 & 1.126 / 1.053 & 1.752 / 1.633 \\
OpenAI-3072 & 0.973 & 1.328 / 1.316 & 1.825 / 2.019 \\
\bottomrule
\end{tabular}

\end{minipage}\par\medskip

两种估计在该有限网格均保留 FP32 数值优势。OpenAI-3072 的 RaBitQ 插值区间为
recall 0.9669--0.9796，揭示较强插值敏感性。这些锚点不替换 E1 完整网格，
不建立通用优势幅度，也不是 FP32 重训。

\section{材料获取}
\label{sec:b-artifacts}
可获得的源码快照、测量记录与分析工具见
\href{https://github.com/lrlorry/8800_report/blob/main/docs/artifact-access.md}{项目复现材料}，
包括 v71 源码／测量包及部分 T4/T5、T1、A2 输入。完整向量、训练状态与全部历史
可执行文件未包含在内；GPU 复现还需要这些输入及匹配的构建环境。
